\section{Impact}
\label{sec:Impact}

The proposed NLU-Solr-GIS architecture provides a reusable, domain-agnostic open-source software pattern for conversational spatial discovery, unlocking new scientific workflows across geospatial data spaces~\cite{softwarex2015}.

\subsection{Cross-Domain Generalization and Architectural Reusability}
\label{subsec:CrossDomain}

A fundamental design objective of this architecture is functional decoupling between intent parsing, domain terminology normalization, spatial indexing, and cartographic rendering. This separation allows developers to transfer the complete conversational GIS pipeline to alternate spatial domains by modifying only two declarative configuration files, without altering core LLM orchestration or GIS synchronization logic:
\begin{enumerate}
    \item \textbf{OpenAPI Schema Specification}: Redefining domain-specific entity slots and categorical enumerations.
    \item \textbf{Thesaurus and Solr Mapping}: Binding canonical entities to underlying spatial index fields.
\end{enumerate}

To illustrate this adaptability, consider a transfer to \textbf{Cadastral Management and Urban Planning}, where an urban planner searches for available municipal or industrial land:
\begin{quote}
\textit{``Show unbuilt rustic parcels over 5 hectares in Madrid province that are not environmentally protected''}
\end{quote}
Adapting the pipeline requires configuring the OpenAPI schema with target cadastral dimensions—such as parcel zoning classification, minimum land surface area, administrative provinces (e.g., Madrid), and environmental protection status—which are deterministically mapped to the spatial index filter queries. All downstream components—greedy schema decoding, dynamic active filter badges, automatic bounding box re-centering, and animated glowing marker feedback—function immediately out of the box. Identical modular adaptation applies to disaster response (wildfire and flood monitoring) or renewable energy planning (identifying solar farm parcels based on terrain slope and power grid proximity).

\subsection{Scientific Impact and Sovereign Open Science}
\label{subsec:ScientificImpact}

Deployed as the conversational demonstrator for the European project \textbf{CRMsDataSpace} (Grant Agreement No.~101216677), the architecture directly addresses the policy priorities of the European Critical Raw Materials Act (Regulation EU 2024/1252)~\cite{european2024crma}. It enables mining engineers, geologists, and environmental auditors to cross-examine historical extractive waste inventories, identifying secondary critical raw material recovery opportunities while assessing acid mine drainage risks across borders~\cite{vidal2017environmental, lottermoser2010mine}.

Crucially, the architecture eliminates the primary barrier to scientific reproducibility in AI-assisted software:
\begin{itemize}
    \item \textbf{Complete Data Sovereignty}: Native execution of open-weight foundation models on institutional private GPU clusters guarantees that sensitive geospatial and industrial records remain strictly on-premises, fully compliant with European Data Spaces regulations and GAIA-X principles.
    \item \textbf{Zero-Setup Standalone Reproducibility}: The standalone mock execution mode provides immediate reviewer verification using standard Python with zero external API key requirements or third-party cloud costs.
    \item \textbf{Open Science Availability}: The codebase, complete synthetic 100-site European dataset, and evaluation suite are made fully accessible under the permissive MIT License in a public GitHub repository.
\end{itemize}
